% chapters/ch27_time_horizon.tex
% Chapter 27: TimeHorizonLayer — v7.2 Revision
% Supersedes: v7.1 display-only TimeHorizonLayer (INV-2.1-6 regime)
% Source of revision: v7.2 foundational discussion, 2026-08-27.
% Cross-reference: vNext simulator horizon profiles (Part VIII)
%                  multi-horizon validation (Part XII §1.3 Performance).

\chapter{TimeHorizonLayer (v7.2: Horizon-Integrated Core)}
\label{ch:time-horizon-layer}

% =====================================================================
\section*{v7.2 Revision Note}
\addcontentsline{toc}{section}{v7.2 Revision Note}

\nrmobox{Root-cause finding (v7.1 retrospective)}{
The v7.1 TimeHorizonLayer was architecturally inert: it computed a
multi-horizon \texttt{warn\_flag} but was explicitly forbidden
(INV-2.1-6) from influencing \texttt{NRMO\_CORE.veto()}. This kept
Core's single-judge property (INV-4) intact, but at the cost of
making the layer diagnostically visible yet causally powerless. A
horizon-aware warning that cannot alter the verdict does not prevent
the failure it detects.
}

\paragraph{v7.2 resolution.}
Rather than retaining an external second layer, TimeHorizonLayer is
now integrated \emph{inside} \texttt{NRMO\_CORE.veto()} as a native
multi-horizon evaluation. There remains exactly one judge (INV-4 is
preserved by construction, not by exclusion), but that judge now
evaluates candidate actions across five horizons before rendering a
single verdict.

% =====================================================================
\section{Two Ruin Concepts, Clarified}
\label{sec:time-horizon-two-ruins}

The v7.2 revision required first resolving a conflation present since
early drafts: True Ruin and Passive Ruin were treated as parallel
failure categories. They are not. They are two diagnostic points in a
single causal chain.

\nrmobox{True Ruin (unchanged, absolute)}{
True Ruin remains a single, absolute, horizon-independent condition:
a state from which no recovering trajectory exists, under \emph{any}
horizon. Its definition is not relativized by this revision — doing
so would erode the boundary A1--A3 exists to protect.
}

\nrmobox{Passive Ruin (redefined)}{
Passive Ruin is \textbf{not} the absence of action. It is the
selection of an action, while an avoidability window is still open,
that irreversibly narrows or closes that window — typically without
the chooser recognising the action in those terms. ``Passive'' refers
to the chooser's relationship to the significance of the choice, not
to the presence or absence of motion.
}

\paragraph{Ordering.} True Ruin is downstream (a confirmed state);
Passive Ruin is upstream (a choice-time event). This ordering is the
basis for the v7.2 architectural change: a downstream-only detector
(v7.1) can only warn after the choice that mattered has already been
made. An upstream detector can intervene before execution.

% =====================================================================
\section{Avoidability Window}
\label{sec:time-horizon-window}

\begin{lstlisting}[basicstyle=\ttfamily\small,frame=single]
def avoidability_window(state, horizon):
    """
    Remaining time/optionality, at this horizon, before the
    trajectory toward True Ruin becomes locked in.
    """
    trajectory = project_to_true_ruin(state, horizon)
    return time_or_optionality_remaining_before_lock_in(trajectory)
\end{lstlisting}

The measurement function for \texttt{avoidability\_window} is left
\textbf{unspecified in this chapter}. This is a known open item
(Section~\ref{sec:time-horizon-open-items}), not an oversight: a
premature numeric definition here would misrepresent the state of
the theory.

% =====================================================================
\section{Horizon-Integrated Core}
\label{sec:time-horizon-core-integration}

\begin{lstlisting}[basicstyle=\ttfamily\small,frame=single]
def NRMO_CORE.veto(state, candidate_action, horizon_set):
    verdicts = {}
    for horizon in horizon_set:
        window_before = avoidability_window(state, horizon)
        projected = apply(state, candidate_action)

        if is_true_ruin(projected, horizon):
            verdicts[horizon] = "REJECT_TRUE_RUIN"
            continue

        window_after = avoidability_window(projected, horizon)
        if window_before.is_open and window_after.triggers_lock_in():
            verdicts[horizon] = "REJECT_PASSIVE_RUIN"
            continue

        verdicts[horizon] = "ALLOW"

    final = most_restrictive(verdicts)
    return {
        "verdict": final,
        "horizon_breakdown": verdicts,
        "binding_horizon": binding(verdicts),
    }
\end{lstlisting}

\paragraph{Preservation of INV-4.} There is still exactly one
verdict-producing function. INV-4 (Core's sole admissibility
authority) is preserved because horizon-awareness is now internal to
Core, not a competing external source. INV-2.1-6, which prohibited
wiring an external \texttt{warn\_flag} into Core, becomes moot under
this design: the entity it was written to keep out no longer exists
as a separate entity. INV-2.1-6 is retained in the historical record
(Chapter~27, v7.1 text, archived) rather than deleted, per the
project's practice of not silently erasing superseded invariants.

\paragraph{Asymmetry between the two verdict types.} True Ruin
rejection is unconditional Minimax across horizons: one positive
detection anywhere is sufficient to REJECT. Passive Ruin rejection
depends on the window/action relationship and is comparatively
softer — this asymmetry is intentional and mirrors the conceptual
asymmetry established in Section~\ref{sec:time-horizon-two-ruins}.

% =====================================================================
\section{Execution Flow Change}
\label{sec:time-horizon-flow-change}

The v2.1 Integrated Execution Flow (v7.1, Chapter 13) placed
TimeHorizonLayer at Step 5, \emph{after} mode selection and
\emph{before} the FROZEN v2.0 pipeline — evaluating a decision that
was, in effect, already made. v7.2 requires evaluation
\emph{before} an action is committed:

\begin{lstlisting}[basicstyle=\ttfamily\small,frame=single]
v7.2 flow (candidate-evaluation form):
1-4. [unchanged: ReflexLayer, ModeSelector, NonErgodicMonitor,
      MetaGovernanceLayer pre_check]
5'. Candidate action(s) generated (not yet executed)
6'. NRMO_CORE.veto(state, candidate_action, horizon_set) per candidate
7'. Candidates verdicted REJECT_TRUE_RUIN / REJECT_PASSIVE_RUIN excluded
8'. Execute from remaining ALLOW candidates
\end{lstlisting}

\nrmobox{Scope of this change}{
This is a structural change to a pipeline whose v2.0 core was
FROZEN and whose v7.1 integration was validated across three
simulation horizons (200/500/1000 steps). It is \textbf{not}
validated at time of writing. Given the v8.x precedent (Stagnation
world regression to 33\% Pareto pass rate, subsequently rolled
back), this chapter is published as a \textbf{proposal-grade
specification}, not a validated release component. See
Section~\ref{sec:time-horizon-open-items}.
}

% =====================================================================
\section{Open Items}
\label{sec:time-horizon-open-items}

\begin{enumerate}
\item \textbf{Avoidability window measurement} is undefined. This is
the largest blocking gap between this specification and an
implementation.
\item \textbf{Candidate-action generation}: the v7.1 pipeline does
not clearly generate multiple explicit candidates upstream of Core;
this may require changes outside this chapter's scope.
\item \textbf{Lock-in threshold}: the boundary between ``window
temporarily narrowed, still recoverable'' and ``window irreversibly
closed'' is not yet formalised.
\item \textbf{StrongEngine symmetry}: the forward-progress engine
requires an analogous (inverted) treatment — actions that
\emph{widen} the window are the object of optimisation there, rather
than actions to be filtered out.
\item \textbf{Validation}: no simulation results exist yet for this
design. Given the project's v8.x history, staged validation
(simulation-only, before any live/operational use) is a
precondition, not an afterthought.
\end{enumerate}

% =====================================================================
\section{Relationship to v7.1 Text}
\label{sec:time-horizon-v71-relationship}

The v7.1 TimeHorizonLayer specification (display-only, five fixed
horizons, INV-2.1-6) is preserved verbatim in
\texttt{archive/ch27\_time\_horizon\_v7\_1.tex} for historical
continuity. This chapter does not erase that text; it supersedes it
for v7.2 onward while keeping the prior design inspectable.


% =====================================================================
\section{Implementation Mapping (v7.2)}
\label{sec:time-horizon-implementation-mapping}

\nrmobox{Reused, not rebuilt}{
The Monte Carlo rollout-survival logic already present in
\texttt{engine/omega\_full.py::score\_candidate()} (the
\texttt{survived\_ratio} computation) is structurally equivalent to
the Viability Kernel membership approximation this chapter requires.
Rather than reimplementing it, the rollout mechanism was
\emph{extracted} into a governance-callable shared module.
}

\paragraph{New files (v7.2, additive only).}
\begin{itemize}
\item \texttt{core/viability.py} — \texttt{viability\_score()},
\texttt{window\_open()}, \texttt{closes\_window()},
\texttt{passive\_ruin\_window\_signal()}. Deliberately excludes
engine-specific scoring (drift penalty, portfolio synergy, failure
memory) to preserve the governance/execution boundary.
\item \texttt{governance/nrmo\_core.py} — \texttt{nrmo\_v72\_veto()}
and \texttt{construct\_admissible\_set\_v72()}, added alongside the
existing \texttt{nrmo\_veto()} / \texttt{nrmo\_vnext\_veto()} /
\texttt{construct\_admissible\_set()}, which remain unmodified and
remain the default path. v7.2 is opt-in, not wired into any existing
call site by this revision.
\end{itemize}

\paragraph{Legacy Passive Ruin detector retained as lagging indicator.}
\texttt{core/ruin.py::update\_passive\_ruin()} (streak-based:
consecutive-step thresholds on $O$, $K$, and compound decline)
predates this chapter's choice-time redefinition of Passive Ruin
(Section~\ref{sec:time-horizon-two-ruins}). It is \textbf{not}
removed. It is retained as a lagging/backstop signal, complementing
the new leading (choice-time, window-closure) signal — see
\texttt{docs/v72\_implementation\_integration\_map.md} for the
two-tier detection design.

\paragraph{Verification path (reusing existing infrastructure).}
Chapter 20's real-side/civ-sim two-layer structure (Strong Engine
$\Omega$ Full vs.\ MaxForwardEngine) provides a validation pathway
that did not need to be built from scratch: v7.2's Core extension is
to be trialled on the civ-sim side (MaxForwardEngine) against the
existing v7.1 (single-horizon, instant-threshold) veto before any
consideration of real-side deployment. This mirrors the project's
post-v8.x posture of staged, simulation-first validation.

\paragraph{Status.} All v7.2 code is proposal-grade and
\textbf{unvalidated} — no simulation runs have been performed against
it at time of writing. See Open Items
(Section~\ref{sec:time-horizon-open-items}).
